% !TeX root = ../../Book1.tex
% 纲要标题：### E.1 域上的序与形式实性
% 纲要位置：计划纲要.md，第 914 行。
% BEGIN APPENDIX OUTLINE SYNC
% #### E.1.1 有序域、正锥与平方和
%
% #### E.1.2 形式实域与 Artin–Schreier 判据
%
% #### E.1.3 域上的序与二次型
% END APPENDIX OUTLINE SYNC

\section{域上的序与形式实性}
\label{b1:sec:E01}

说一个数为正，需要与加法、乘法相容的比较规则。域公理本身不含这样的规则：复数域不能带域序，而同一个有理函数域却可以有许多不同的序。本节先给出正锥的代数条件，再将可赋序性转成平方和问题。

% 纲要内容（仅作写作计划，尚非教材正文）：
%
% * 有序域、正锥与平方和；实数域和有理函数域的例子
% * 形式实域的定义：\(-1\) 不是有限个平方之和；Artin–Schreier 可赋序判据
% * 指定一个序与枚举域上的全部序是不同问题；与第二卷二次型的联系
%

\subsection{有序域、正锥与平方和}
\label{b1:subsec:E01:01}

\begin{definition}\label{b1:def:ordered-field}
\term[youxu-yu]{有序域}[ordered field]是域 \(K\) 连同全序 \(\leq\)，使 \(a\leq b\) 蕴含 \(a+c\leq b+c\)，并使 \(a,b\geq0\) 蕴含 \(ab\geq0\)。它的\term[zheng-zhui]{正锥}[positive cone]在本附录采用含零约定，即 \(P=\{a\in K\mid a\geq0\}\)。
\end{definition}

正锥的等价描述为
\[
P+P\subseteq P,\qquad PP\subseteq P,\qquad
P\cup(-P)=K,\qquad P\cap(-P)=\{0\}.
\]
由这些条件反过来定义 \(a\leq b\) 当且仅当 \(b-a\in P\)，可分别检验自反性、传递性、反对称性、全序性与运算相容性。非零平方总为正：非零 \(a\) 与 \(-a\) 必有一个为正，而两者的平方相同。特别 \(1>0\)，所以所有正整数的像为正；有序域必有特征零。

\ProseParagraphBreak
通常的 \(\mathbb Q\) 和 \(\mathbb R\) 是有序域。对 \(\mathbb Q(t)\)，规定非零有理函数 \(f/g\) 为正当且仅当 \(f,g\) 的最高次项系数同号。这与表示无关，因为换表示时交叉乘积相等；乘法保持正号，加法可先通分到正的平方分母，两个正最高次项不会相消。因此它给出一个序，且 \(t>n\) 对所有整数 \(n\) 成立。以 \(t\mapsto-t\) 拉回这个序，则得到 \(-t>n\) 的另一个序。

\begin{definition}\label{b1:def:sum-of-squares}
域 \(K\) 的平方和集合记为
\[
\Sigma K^2=\left\{\sum_{j=1}^r a_j^2\ \middle|\ r\geq0,\ a_j\in K\right\};
\]
空和为零。若 \(-1\notin\Sigma K^2\)，称 \(K\) 为\term[xingshi-shiyu]{形式实域}[formally real field]。
\end{definition}
\symidx{\(\Sigma K^2\)：域中有限个平方的和}

任一域序都使平方和非负，所以可赋序的域必形式实。形式实性只涉及四则运算，没有预先选择哪个元素为正。有限域不满足它，因为特征 \(p\) 时 \(-1\) 是 \(p-1\) 个 \(1^2\) 的和；\(\mathbb C\) 中则已有 \(-1=i^2\)。

\subsection{形式实域与 Artin–Schreier 判据}
\label{b1:subsec:E01:02}

下面的\term[Artin-Schreier-kafuxu-panju]{Artin--Schreier 可赋序判据}[Artin--Schreier orderability criterion]把上一段的必要条件变成充分条件。它与\secref{b1:sec:2004}中正特征循环扩张的 Artin--Schreier 理论属于不同的问题。

\begin{theorem}[Artin--Schreier]\label{b1:thm:artin-schreier-ordering}
域 \(K\) 能够赋予一个域序，当且仅当它形式实。
\end{theorem}
\begin{proof}
必要性已经说明。设 \(K\) 形式实。考虑所有子集 \(T\subseteq K\)，要求它含有全部平方，对加法、乘法封闭，且不含 \(-1\)。这一集合族非空，因为 \(\Sigma K^2\) 满足条件。按包含关系取任一链，其并仍含全部平方且不含 \(-1\)；加、乘涉及有限个元素，能在链中同一个成员里检验，所以并仍满足条件。由\cref{b1:thm:A-zorn}，存在极大成员 \(P\)。

若 \(s\in P\) 且 \(s\neq0\)，则 \(s^{-1}=s(s^{-1})^2\in P\)。所以 \(a,-a\in P\) 且 \(a\neq0\) 将推出 \(-1=(-a)a^{-1}\in P\)，不可能。因此 \(P\cap(-P)=\{0\}\)。

还需证明每个 \(a\) 或其负元属于 \(P\)。若 \(a\notin P\)，考察
\[
T=P-aP=\{r-as\mid r,s\in P\}.
\]
它包含 \(P\)，对加法封闭；乘法展开中的 \(a^2\) 属于 \(P\)，故也对乘法封闭。若 \(-1=r-as\)，则 \(s\neq0\)，从而 \(a=(1+r)s^{-1}\in P\)，矛盾。所以 \(-1\notin T\)，极大性迫使 \(T=P\)。取 \(r=0,s=1\) 得 \(-a\in P\)。至此正锥的全部条件成立，\cref{b1:def:ordered-field}后的构造给出所需域序。
\end{proof}

这个证明还说明：若一个集合已经含全部平方、对加乘封闭且不含 \(-1\)，就能扩成正锥。这里的极大性只用来完成符号选择，不是计算某个具体序的有限算法。

\subsection{域上的序与二次型}
\label{b1:subsec:E01:03}

指定一个序以后，所有非零平方都为正，但非平方元的符号可能随序改变。例如 \(\mathbb Q(\sqrt2)\) 到 \(\mathbb R\) 的两个嵌入把 \(\sqrt2\) 分别送到正根、负根，拉回两个不同的序。可赋序判据只保证至少存在一个序，不枚举所有序。

\ProseParagraphBreak
平方和也把这一问题连到二次型。对角二次型 \(q(x_1,\ldots,x_n)=a_1x_1^2+\cdots+a_nx_n^2\) 的系数在某个固定序中全为正时，任何非零向量都使 \(q\) 为正。例如 \(x^2+y^2=0\) 在有序域中只能有零解；若允许 \(-1\) 为平方，结论立即失败。更一般地，形式实性等价于任意平方和为零时各项均为零，因为只要有一个非零项，就可移项并除以它的平方得到 \(-1\) 的平方和表示。第二卷对二次型作系统研究时，还要区分等价、各向同性与所选域序；本节没有给出一般二次型的分类。
